\documentclass[12pt]{article}
\usepackage[margin=1in]{geometry}
\usepackage{amsmath}
\usepackage{amssymb}
\usepackage{enumitem}

\setlist{topsep=4pt,parsep=2pt}

\title{Discrete Mathematics --- Worksheet 4: Induction and Counting}
\author{FlashTeX Course Staff}
\date{\today}

\begin{document}
\maketitle

\noindent Name: \hrulefill \qquad Section: \hrulefill

\section*{Part A: Warm-up (arabic labels)}

\begin{enumerate}[label=\arabic*.]
  \item Compute $\sum_{k=1}^{5} k^2$.

  \vspace{1cm}
  \hrulefill

  \item Solve for $x$: $2x + 7 = 19$.

  \vspace{1cm}
  \hrulefill

  \item Show that $n^2 - n$ is even for every integer $n$.

  \vspace{1.5cm}
  \hrulefill
\end{enumerate}

\section*{Part B: Nested Structure}

Work through the following multi-part induction problem. Each top-level
item has lettered sub-parts, and one sub-part has roman-numeral sub-sub-parts.

\begin{enumerate}
  \item Let $P(n)$ be the statement $1 + 2 + \cdots + n = \dfrac{n(n+1)}{2}$.
  \begin{enumerate}[label=(\alph*)]
    \item Verify $P(1)$.

    \vspace{0.5cm}

    \item State the inductive hypothesis $P(k)$.

    \vspace{0.5cm}

    \item Prove $P(k) \implies P(k+1)$ by completing each sub-step below.
    \begin{enumerate}[label=\roman*.,itemsep=0pt]
      \item Write down $P(k+1)$ explicitly.
      \item Substitute the inductive hypothesis for the first $k$ terms.
      \item Simplify the resulting expression to match the closed form.
    \end{enumerate}

    \vspace{1cm}
    \hrulefill
  \end{enumerate}

  \item Let $Q(n)$ be the statement that $3 \mid (4^n - 1)$.
  \begin{enumerate}[label=(\alph*)]
    \item Verify $Q(0)$.
    \item Prove $Q(k) \implies Q(k+1)$, showing your algebra with the
          display below:
          \[
            4^{k+1} - 1 = 4 \cdot 4^k - 1 = 4(4^k - 1) + 3.
          \]
          Explain why this proves the divisibility step.

    \vspace{1cm}
    \hrulefill
  \end{enumerate}
\end{enumerate}

\section*{Part C: Compact Lists}

\begin{enumerate}[leftmargin=*]
  \item State the Pigeonhole Principle in your own words.
  \item Give one real-world example where it applies.
\end{enumerate}

The following short list should have no extra vertical spacing between
items:

\begin{enumerate}[nosep]
  \item $2 + 2 = 4$
  \item $3 \times 3 = 9$
  \item $\sqrt{16} = 4$
\end{enumerate}

\section*{Part D: Terminology}

\begin{description}[style=nextline]
  \item[Bijection] A function that is both injective and surjective,
        establishing a one-to-one correspondence between two sets.
  \item[Pigeonhole Principle] If $n$ items are placed into $m$ containers
        with $n > m$, then at least one container holds more than one item.
  \item[Well-Ordering Principle] Every nonempty set of nonnegative integers
        has a least element.
\end{description}

\section*{Part E: Checklist}

\begin{itemize}[label=$\square$]
  \item I verified the base case explicitly.
  \item I stated the inductive hypothesis clearly.
  \item I showed the inductive step follows logically.
  \item I double-checked my algebra in the displayed equations.
\end{itemize}

\section*{Part F: Challenge Problem}

Prove, by induction on $n \geq 1$, that the sum of the interior angles of a
convex polygon with $n+2$ sides is $180n$ degrees. Use the space below,
including at least one displayed equation such as
\[
  S(n+1) = S(n) + 180,
\]
to justify the inductive step.

\vspace{2cm}
\hrulefill

\vspace{0.3cm}
\hrulefill

\end{document}
